\documentclass[11pt,a4paper]{article}
\usepackage[utf8]{inputenc}
\usepackage[T1]{fontenc}
\usepackage[italian]{babel}
\usepackage[margin=2.5cm]{geometry}
\usepackage{xcolor}
\usepackage{enumitem}
\usepackage{hyperref}
\usepackage{fancyhdr}
\usepackage{titling}
\usepackage{tabularx}
\usepackage{booktabs}
\usepackage{textcomp}

\definecolor{studioteal}{HTML}{0D6E6E}
\definecolor{studiogrigio}{HTML}{5B6B6A}

\hypersetup{
  colorlinks=true,
  linkcolor=studioteal,
  urlcolor=studioteal,
  citecolor=studioteal,
  pdftitle={Guida per i pazienti},
  pdfauthor={Studio Medico Dottoressa Braglia}
}

\pagestyle{fancy}
\fancyhf{}
\fancyhead[L]{\small\color{studiogrigio} Studio Medico Dottoressa Braglia}
\fancyhead[R]{\small\color{studiogrigio} Guida per i pazienti}
\fancyfoot[C]{\small\thepage}
\renewcommand{\headrulewidth}{0.4pt}
\renewcommand{\headrule}{\hbox to\headwidth{\color{studioteal}\leaders\hrule height \headrulewidth\hfill}}

\setlist[itemize]{leftmargin=1.4em, itemsep=0.25em}
\setlist[enumerate]{leftmargin=1.6em, itemsep=0.35em}

\newcommand{\schermata}[1]{\textbf{\color{studioteal}#1}}
\newcommand{\bottone}[1]{\textbf{``#1''}}

\title{\color{studioteal}\Huge Guida per i pazienti}
\author{\Large Studio Medico Dottoressa Braglia \\[.3em] \normalsize Ambulatori di Arceto e Casalgrande}
\date{}

\begin{document}

\begin{titlepage}
\centering
\vspace*{3cm}
{\color{studioteal}\Huge\bfseries Guida per i pazienti}\\[1em]
{\Large Come usare il sito dello studio}\\[2em]
{\large Studio Medico Dottoressa Braglia}\\
{\large Ambulatori di Arceto e Casalgrande}\\[3em]
{\normalsize\color{studiogrigio} \url{https://studiomedicobragliavaleria.it}}
\vfill
{\small\color{studiogrigio} Questa guida spiega come creare un account, prenotare una visita, \\
richiedere medicinali o esami, e gestire il proprio account dal sito.}
\end{titlepage}

\tableofcontents
\newpage

\section{Come accedere al sito}

Il sito è raggiungibile all'indirizzo \url{https://studiomedicobragliavaleria.it}, da computer, tablet o telefono. Funziona con qualsiasi browser aggiornato (Chrome, Safari, Edge, Firefox) ed è sempre disponibile, giorno e notte.

La prima cosa che si vede è la schermata di accesso: da lì si entra con un account già esistente, oppure se ne crea uno nuovo.

\section{Creare un account}

\begin{enumerate}
  \item Nella schermata di accesso, cliccare su \bottone{Crea account}.
  \item Compilare il modulo con \textbf{nome}, \textbf{cognome}, \textbf{telefono}, \textbf{email} e una \textbf{password} di almeno 10 caratteri.
  \item Cliccare su \bottone{Registrati}.
\end{enumerate}

A quel punto il sito manda un'email all'indirizzo indicato, con un link per confermarlo. \textbf{Senza aprire quel link non si può accedere.} Il link resta valido 24 ore.

\begin{itemize}
  \item Se l'email non arriva entro un minuto, controllare nella cartella \textbf{spam / posta indesiderata}.
  \item Se il link scade prima di essere aperto, dalla pagina di accesso si può cliccare su \bottone{Rinvia il link} per riceverne uno nuovo.
\end{itemize}

\subsection*{Perché a volte arriva ``Hai già un account''}

Se si prova a registrarsi con un indirizzo già usato in passato, il sito non lo dice apertamente (per non rivelare a chiunque quali email sono già nostre pazienti): manda comunque un'email, ma il suo contenuto sarà diverso — dirà che quell'indirizzo ha già un account, invitando ad accedere con la password, oppure a reimpostarla se dimenticata.

\section{Accedere}

Dalla schermata iniziale, inserire l'email e la password usate in fase di registrazione, poi cliccare su \bottone{Accedi}.

\subsection{Password dimenticata}

Se non si ricorda la password:
\begin{enumerate}
  \item Nella schermata di accesso cliccare su \bottone{Reimpostala}.
  \item Scrivere l'email con cui ci si è registrati e inviare.
  \item Se l'indirizzo corrisponde a un account, arriva un'email con un link per scegliere una password nuova. Il link resta valido \textbf{2 ore}.
\end{enumerate}

\subsection{Email dimenticata}

Se invece non si ricorda \emph{quale} email è stata usata per la registrazione:
\begin{enumerate}
  \item Nella schermata di accesso cliccare su \bottone{Recuperala}.
  \item Scrivere il numero di \textbf{telefono} usato in fase di registrazione.
  \item Se corrisponde a un account, arriva un promemoria dell'indirizzo email — \textbf{alla casella stessa}, mai mostrato sullo schermo, per sicurezza.
\end{enumerate}

Se non si ricorda nemmeno il telefono usato, l'unica strada è contattare direttamente lo studio: il personale può verificare i dati e correggerli dal pannello.

\section{Prenotare una visita}

\begin{enumerate}
  \item Dal menu in alto, aprire \schermata{Prenota}.
  \item Scegliere l'\textbf{ambulatorio}: \textbf{si può prenotare solo in quello di residenza.}
  \item Scegliere un giorno dal calendario: i giorni con il pallino verde hanno posti liberi.
  \item Scegliere un orario libero.
  \item Scrivere il motivo della visita e confermare.
\end{enumerate}

Arriva subito un'email di conferma con il \textbf{codice della prenotazione}: conservarlo, serve per consultarla o annullarla in seguito.

\textbf{Nota:} le visite richieste dal sito nascono come \emph{richieste}: lo studio le conferma (o le rifiuta, indicandone il motivo) entro breve, e a quel punto arriva una seconda email con la conferma definitiva.

\section{Richiedere medicinali, una visita specialistica o esami del sangue}

Dal menu in alto sono disponibili tre moduli distinti:

\begin{itemize}
  \item \schermata{Medicinali} --- per chiedere il rinnovo di una ricetta.
  \item \schermata{Visite specialistiche} --- per chiedere l'impegnativa di una visita da uno specialista.
  \item \schermata{Esami del sangue} --- per chiedere l'impegnativa degli esami.
\end{itemize}

In ciascun modulo si indicano nome, cognome, telefono, email e cosa serve. Per le visite specialistiche e gli esami è possibile \textbf{allegare la foto} di una richiesta già scritta da uno specialista: da telefono si può scattarla direttamente.

Dopo l'invio arriva un'email di conferma con un codice. Quando lo studio risponde (approvando, magari con qualche modifica concordata, oppure spiegando perché non può essere accolta) arriva una seconda email.

\section{Controllare o annullare}

Nella sezione \schermata{La mia richiesta} si inserisce il codice ricevuto via email per vedere lo stato di una prenotazione o di una richiesta, ed eventualmente annullarla:

\begin{center}
\begin{tabularx}{\linewidth}{lX}
\toprule
\textbf{Prefisso del codice} & \textbf{Cosa indica} \\
\midrule
\texttt{PRE-\dots} & Prenotazione di una visita \\
\texttt{MED-\dots} & Richiesta di medicinali \\
\texttt{SPE-\dots} & Richiesta di visita specialistica \\
\texttt{ESA-\dots} & Richiesta di esami del sangue \\
\bottomrule
\end{tabularx}
\end{center}

Chi ha effettuato l'accesso trova anche le proprie pratiche direttamente in quella sezione, senza dover ricordare i codici a memoria.

\section{Il proprio account}

Una volta effettuato l'accesso, in alto a destra si trova il menu \schermata{Account}. Da lì, \bottone{Modifica account} permette di cambiare:

\begin{itemize}
  \item nome, cognome e telefono, liberamente;
  \item l'email con cui si accede --- serve però la \textbf{password attuale}, come conferma;
  \item la password --- serve la password attuale e una nuova di almeno 10 caratteri.
\end{itemize}

\section{Notifiche sul telefono}

Il sito può mandare un avviso direttamente sul telefono o sul computer (oltre alle email) quando una prenotazione viene confermata, una richiesta viene evasa, e così via.

Per attivarle: menu in alto, pulsante \bottone{Attiva notifiche}. Il browser chiederà un permesso: basta accettarlo.

\subsection*{Nota per iPhone e iPad}

Su iPhone e iPad, Apple permette le notifiche solo ai siti \textbf{aggiunti alla schermata Home} --- vale per qualsiasi browser si usi, anche Chrome. Per attivarle:
\begin{enumerate}
  \item Aprire il sito con \textbf{Safari}.
  \item Toccare l'icona \textbf{Condividi}.
  \item Scegliere \textbf{``Aggiungi alla schermata Home''}.
  \item Aprire l'icona appena creata: da lì il pulsante ``Attiva notifiche'' funziona.
\end{enumerate}

\section{Installare il sito come app}

Il sito si può installare sul telefono o sul computer come una vera app, gratuitamente, senza passare da nessun negozio di app.

\begin{itemize}
  \item \textbf{Android, Chrome, computer}: cercare il pulsante \bottone{Installa l'app} in alto nella pagina.
  \item \textbf{iPhone / iPad}: seguire gli stessi passaggi descritti sopra per le notifiche (Safari \textrightarrow\ Condividi \textrightarrow\ Aggiungi alla schermata Home): è lo stesso identico gesto.
\end{itemize}

\section{L'assistente virtuale}

In basso a destra sul sito è disponibile un piccolo assistente (chatbot) che guida passo dopo passo nella prenotazione di una visita, in una richiesta di medicinali/specialistica/esami, oppure nel controllo di una pratica già aperta con il suo codice.

Alcuni comandi utili durante la conversazione:
\begin{itemize}
  \item scrivere \textbf{``indietro''} per correggere l'ultima risposta, senza ricominciare tutto;
  \item scrivere \textbf{``menu''} per tornare al punto di partenza;
  \item scrivere \textbf{``aiuto''} per un riepilogo di cosa sa fare, incluso come cambiare password, dati dell'account o attivare le notifiche.
\end{itemize}

\section{I nostri ambulatori}

\begin{center}
\begin{tabularx}{\linewidth}{lX}
\toprule
\multicolumn{2}{l}{\textbf{Ambulatorio di Arceto}} \\
\midrule
Indirizzo & Via Piazza Castello, 10 --- 42019 Arceto (RE) \\
Telefono & 329 154 5236 \\
Orari & Lunedì, martedì, giovedì, venerdì: 11:00--13:00 \\
       & Mercoledì: 17:30--19:00 \\
\bottomrule
\end{tabularx}

\vspace{1em}

\begin{tabularx}{\linewidth}{lX}
\toprule
\multicolumn{2}{l}{\textbf{Ambulatorio di Casalgrande}} \\
\midrule
Indirizzo & Via Canale, 29 --- 42013 Casalgrande (RE) \\
Telefono & 347 245 0118 \\
Orari & Lunedì, giovedì: 17:30--19:00 \\
       & Mercoledì, venerdì: 08:30--10:00 \\
\bottomrule
\end{tabularx}
\end{center}

\textbf{Si ricorda di prenotare solo nell'ambulatorio di residenza.}

\section{In caso di emergenza}

Per un'emergenza medica chiamare sempre il \textbf{112}. Per urgenze non differibili al di fuori dell'orario di apertura, contattare la Guardia Medica.

\end{document}
